% Chapter 4 - Development, testing and results.
% Largest chapter. Sources: docs/requirements.md, docs/architecture.md,
% docs/diagrams.md, claude/decisiones.md (DEC-001 to DEC-036), code
% under src/, tests under tests/.
% Target length: 18-22 pages.

\chapter{Development, testing and results}
\label{ch:desarrollo}

This chapter is the technical core of the report. It walks through
the development cycle of the proof of concept in five blocks:
requirements analysis, architectural design, implementation, testing
strategy and functional demonstration. It closes with a qualitative
comparison against the two libraries the supervisor had written
before this work. The detailed requirements catalogue is moved to
Appendix~\ref{ap:requirements}; the version presented here is the
structured summary that supports the rest of the chapter.

%==================================================================
\section{Requirements analysis}
\label{sec:analisis}

The requirements analysis was carried out as Phase~2 of the project
and produced \texttt{docs/requirements.md} in the repository. The
analysis follows a lightweight format suited to a proof of
concept: actors and use cases in natural language, functional
requirements grouped by area, a short table of non-functional
requirements, and an explicit list of items that are out of scope
for v1. We summarize the structure here; the full catalogue is in
Appendix~\ref{ap:requirements}.

\subsection{Actors}

\Cref{tab:actors} lists the actors of the system. The role design
follows DEC-011 and DEC-021: three roles with cascade inheritance
through the tree, and a single user that can hold different roles in
different branches.

\begin{table}[ht]
  \centering
  \small
  \begin{tabular}{p{0.22\linewidth} p{0.70\linewidth}}
    \toprule
    \textbf{Actor} & \textbf{Description} \\
    \midrule
    Administrator (\emph{superuser}) &
    Manages hypervisors, base templates, users, the full tree and
    global configuration. \\
    Manager &
    A professor or equivalent. Manages assigned branches of the
    tree. Uses both the core and the teaching application. \\
    End user &
    A student or equivalent. Operates only the VMs assigned to them. \\
    Hypervisor &
    External system. UnivOrch talks to it through a connector;
    the hypervisor does not initiate actions. \\
    System &
    UnivOrch itself, running the job engine and automatic tasks
    (backup rotation, retention cleanup). \\
    \bottomrule
  \end{tabular}
  \caption{Actors of UnivOrch.}
  \label{tab:actors}
\end{table}

\subsection{Use cases}

Use cases are split into three groups. \emph{Core} use cases are
generic operations of the orchestrator (authentication, defining
hypervisors and templates, creating folders, deploying and
undeploying VMs, viewing job history). \emph{Teaching application}
use cases (UC-TA-1\dots) operate at the domain level: deploying a
whole subject from a roster of students, regenerating a student's
environment, end-of-term recycling. \emph{System} use cases describe
the automatic behaviour of the daemon (backup, retention cleanup,
detection of jobs left in flight after a crash). The full catalogue
contains 32 use cases.

The most representative use cases of the core are UC-ADM-2 (define a
hypervisor with its credentials and datastores), UC-ADM-3 (define a
base template), UC-MGR-1 (a manager creates folders inside their
branch), UC-MGR-3 (a manager loads a YAML definition of a group of
VMs), UC-EU-1 (an end user operates their own VMs through the
simplified view) and UC-ADM-8 (the administrator force-undeploys a
descriptor stuck in \texttt{broken} state).

\subsection{Functional requirements}

The functional requirements are grouped by area: descriptor tree
and hierarchical organization; declarative load and inheritance;
hypervisor connectors and life-cycle operations; jobs and
operations; persistence and backup; identity, authentication and
authorization; and the teaching application. Within each area, the
requirements stay close to the use cases of the same area and
introduce only the precision needed to bridge into the architecture
phase.

\subsection{Non-functional requirements}

\Cref{tab:nfr} lists the headline non-functional requirements. The
proof of concept is not a benchmark exercise; the numbers below are
soft targets used to keep the implementation honest, not contractual
service levels.

\begin{table}[ht]
  \centering
  \small
  \begin{tabular}{p{0.25\linewidth} p{0.65\linewidth}}
    \toprule
    \textbf{Quality} & \textbf{Target} \\
    \midrule
    Reproducibility &
    A new clone of the repository plus one command must start a
    working development environment. \\
    Portability &
    The same code base must support new hypervisors without changes
    to the core. \\
    Testability &
    The core must be exercisable without a real hypervisor. \\
    Maintainability &
    Every component must have a test, every public function must be
    typed, and the documentation must trace every architectural
    choice back to a numbered decision. \\
    Operability &
    The complete service must run inside a single Docker container
    and be deployable in one command. \\
    \bottomrule
  \end{tabular}
  \caption{Non-functional requirements of UnivOrch v1 (summary).}
  \label{tab:nfr}
\end{table}

\subsection{Out of scope for v1}

The following items are explicitly excluded from this proof of
concept: automatic reconciliation loops, active/passive high
availability, IPAM integration, snapshot management through the
orchestrator, discovery of pre-existing VMs not created by UnivOrch,
multi-tenant isolation beyond what RBAC provides, and any
integration with single-sign-on or directory services.

%==================================================================
\section{Architectural design}
\label{sec:diseno}

The architecture is documented in detail in
\texttt{docs/architecture.md} and visually summarized in
\texttt{docs/diagrams.md}; this section presents the same content in
narrative form, with the diagrams embedded next to the text they
support.

\subsection{Two-layer architecture (DEC-004)}
\label{subsec:two-layers}

UnivOrch is built as two clearly separated layers. The lower layer
is a \emph{generic orchestration engine} that knows about folders,
descriptors, jobs, connectors and inheritance, but knows nothing
about academic subjects, students or any other domain. The upper
layer is a \emph{domain application} that translates the vocabulary
of a specific use case into operations of the engine. The teaching
application of this TFG is the only upper layer implemented;
\Cref{fig:two-layers} also shows where alternative upper layers
would sit (a CTF organizer, a workshop deployer, a research-lab
orchestrator).

The boundary between layers is the public REST API of the daemon.
Every upper layer is a client of this API. The same separation
applies to the operational interfaces (CLI, future web GUI): they
are also clients of the API, not parts of the engine. Nothing in the
engine knows what kind of client is calling it.

\begin{figure}[ht]
  \centering
  \includegraphics[width=0.95\linewidth]{two-layers.pdf}
  \caption{Two-layer overview: the daemon \texttt{univorchd} hosts the
    core; every client (CLI, web GUI, teaching app) speaks REST to it.}
  \label{fig:two-layers}
\end{figure}

\subsection{C4 context and containers}
\label{subsec:c4-l1-l2}

\Cref{fig:c4-context} shows the system at C4 level~1. Three kinds of
human or external actor interact with UnivOrch: the administrator
or teacher who defines the tree and triggers deployments; the
student who operates their own machines; and external integrations
(CI/CD pipelines, GitOps tools) that consume the public REST API.
The orchestrator itself talks to the underlying hypervisors (VMware
vSphere, Proxmox VE), and students access their VMs by IP
\emph{directly}, not through the orchestrator -- this distinction
matters when discussing the failure modes.

\begin{figure}[ht]
  \centering
  \includegraphics[width=0.95\linewidth]{context.pdf}
  \caption{C4 level 1 (Context): people and external systems around
    UnivOrch.}
  \label{fig:c4-context}
\end{figure}

\Cref{fig:c4-containers} zooms into the C4 level~2 view. There are
two clearly delimited deployable units inside the system: the
\emph{daemon} \texttt{univorchd} (a Python process running FastAPI
on Uvicorn) and the \emph{CLI} \texttt{univorch} (a Python program
based on cmd2 that talks to the daemon by HTTP). Both live in the
same package distribution and are installed by the same Docker
image; one runs as a long-lived service, the other as an
on-demand command. The future web GUI fits in this view as a third
container, with the same role as the CLI -- a thin client of the
REST API.

\begin{figure}[ht]
  \centering
  \includegraphics[width=0.95\linewidth]{containers.pdf}
  \caption{C4 level 2 (Containers): independently runnable units. The
    daemon \texttt{univorchd} and the CLI \texttt{univorch} are two
    separate binaries; every client speaks HTTP.}
  \label{fig:c4-containers}
\end{figure}

\subsection{Deployment view}
\label{subsec:deployment}

\Cref{fig:deployment} shows the deployment view of Sprint~3. The
target deployment is a single Linux host with Docker that runs the
\texttt{univorchd} container. The container exposes one HTTP port
(default 8080) for the REST API and mounts a Docker named volume
(\texttt{univorch\_data}) where the TinyDB JSON file lives. Both the
named volume and the port are configurable through environment
variables.

The clients sit on tier~1, on any host that can reach the daemon
port. The hypervisor hosts are tier~3; the daemon talks to them
through the connectors implemented in the engine. In the
mock-driven demo, tier~3 is not used: the mock connector lives
inside the daemon and provides a fictitious hypervisor in memory.

\begin{figure}[ht]
  \centering
  \includegraphics[width=0.95\linewidth]{deployment.pdf}
  \caption{Deployment topology since Sprint 3: clients on tier 1, the
    daemon container plus its named volume on tier 2, and the
    hypervisor hosts on tier 3.}
  \label{fig:deployment}
\end{figure}

\subsection{Components (as-built)}
\label{subsec:components}

\Cref{fig:components} is the C4 level~3 view, restricted to what
was actually built. The diagram uses a legend with two states:
\emph{implemented} (solid green) for everything delivered up to
Sprint~3, and \emph{pending} (dashed grey) for everything that was
designed but is left as future work. The pending blocks include the
web GUI, the real hypervisor connectors (VMware and Proxmox) and
the teaching application.

Inside the daemon, the implemented components are: the REST layer
(\texttt{create\_app}, the FastAPI application, the
\texttt{HttpServiceClient} that lives in the CLI), the
\texttt{OrchestratorService} facade, the lexical-closure resolver
with its generic walker (\texttt{\_find\_resource}), the connection
pool indexed by folder path, the connector type registry
(\texttt{CONNECTOR\_TYPES}), the YAML parser, the job engine with
its Commands, the three TinyDB repositories (folders, descriptors,
jobs), the abstract \texttt{HypervisorConnector} and the
\texttt{MockConnector}. The diagram makes the layered structure
explicit: the REST layer depends only on the facade, the facade
depends on the resolver, the job engine, the pool and the parser,
the job engine depends on the connector abstraction and on the
repositories, and the repositories depend on TinyDB.

\begin{figure}[ht]
  \centering
  \includegraphics[width=0.95\linewidth]{components.pdf}
  \caption{C4 level 3 (Components): modules inside \texttt{univorchd}
    after Sprint 3. Solid green = implemented; dashed grey = designed,
    not yet implemented.}
  \label{fig:components}
\end{figure}

\subsection{Domain model}
\label{subsec:domain-model}

The descriptor tree is implemented as a flat collection of records
indexed by their \emph{materialized path} (DEC-008). A folder at
\texttt{/lab/networks} is one record, the descriptor
\texttt{student01} inside it is another record at
\texttt{/lab/networks/student01}, and so on. Navigating the tree
upward or sideways is then a matter of string operations on the
path; there are no pointers between records to maintain. This
choice keeps the TinyDB schema flat and the repository operations
trivial.

Inheritance of properties follows DEC-026: the rule depends on the
data type of the field. A scalar field (\texttt{cpu},
\texttt{memory\_mb}) is overridden by the child; a list field
(\texttt{imports}) accumulates from parent to child; a dictionary
field (\texttt{hypervisors}, \texttt{vm\_templates}) is recursively
merged. The merge stops at the explicit boundary set by the
\texttt{import:} field of each folder, which decides what names
from the parent are visible inside this folder.

The template lexical closure deserves special mention. When a
template is defined in folder $A$ and imported into folder $B$, any
reference inside the template (\emph{e.g.}\ which hypervisor to
target) is resolved \emph{starting from $A$}, not from $B$. The
template carries its definition environment with it. This is the
behaviour of a lexical closure in a programming language and is the
refinement of DEC-012 applied during Sprint~2.

\begin{sidewaysfigure}
  \centering
  \includegraphics[width=0.95\linewidth]{domain-models.pdf}
  \caption{Domain model and connector contract (C4 code level, view
    1/2). Persisted entities (\texttt{Folder}, \texttt{Descriptor},
    \texttt{Job}), input documents (\texttt{DefinitionDocument} and
    friends) and the abstract \texttt{HypervisorConnector}.}
  \label{fig:domain-models}
\end{sidewaysfigure}

\Cref{fig:domain-models} shows the model classes. There are two
families of types. The first family is the \emph{persisted} types
(\texttt{Folder}, \texttt{Descriptor}, \texttt{Job}); they carry an
absolute path or identifier and live in TinyDB. The second family is
the \emph{document} types (\texttt{DefinitionDocument},
\texttt{FolderDef}, \texttt{DescriptorDef}, \texttt{HypervisorDef},
\texttt{VMTemplateDef}); they are produced by the YAML parser and
travel \emph{into} the \texttt{load} operation, where the engine
turns them into persisted entities. The split is intentional:
documents allow relative paths and aliases for the YAML syntax
(\texttt{use hypervisor:}, \texttt{define templates:}), persisted
entities use canonical absolute paths.

\subsection{Patterns applied}
\label{subsec:patterns}

Several classical design patterns shape the implementation. The
relevant ones are listed below; the citation of each pattern is left
to the references in \texttt{docs/architecture.md} to keep the
prose tight.

\begin{description}
  \item[Facade (DEC-031)] -- the \texttt{OrchestratorService} class
        is the single entry point that every client uses. It
        encapsulates the engine and protects clients from the
        internal repositories, the job engine, the connector pool
        and the resolver. The REST layer talks to the facade only.
  \item[Command pattern (DEC-028)] -- every state-changing operation
        is encapsulated in a Command class (\texttt{DeployCommand},
        \texttt{UndeployCommand}, \texttt{StartCommand}, \dots).
        Commands are validated before being run; if validation
        passes, the job engine executes them and tracks the
        outcome in the job log.
  \item[Repository (DEC-007/030)] -- persistence is hidden behind a
        repository class for each aggregate root (folder, descriptor,
        job). The repositories take and return Pydantic models, not
        TinyDB documents. Replacing TinyDB with MongoDB requires
        rewriting only the repository classes.
  \item[Strategy / Abstract Base Class (DEC-029)] -- the
        hypervisor connector is an abstract base class with a small
        and stable interface. The mock connector implements it
        in memory; the future VMware and Proxmox connectors will
        implement it against their respective APIs.
  \item[Protocol (Sprint 3.2)] -- the boundary between the CLI and
        the daemon is described by the \texttt{OrchestratorAPI}
        Protocol. Two classes implement it structurally: the
        in-process \texttt{OrchestratorService} and the
        \texttt{HttpServiceClient}. Clients (the CLI in particular)
        depend on the Protocol, not on either concrete class.
\end{description}

\Cref{fig:engine-rest} shows how these classes relate. The
\texttt{OrchestratorAPI} Protocol unifies the in-process and the
HTTP paths; the engine sits behind the facade; the Commands depend
on repositories and on the connector pool.

\begin{sidewaysfigure}
  \centering
  \includegraphics[width=0.95\linewidth]{engine-rest.pdf}
  \caption{Engine, service and REST boundary (C4 code level, view
    2/2). The \texttt{OrchestratorAPI} Protocol unifies the in-process
    service and the HTTP client; clients depend on the Protocol, not
    on either concrete class.}
  \label{fig:engine-rest}
\end{sidewaysfigure}

\subsection{Declarative load flow (DEC-027)}
\label{subsec:load-flow}

The system follows a strict declarative model. The user does not
issue create or delete commands; they describe what they want by
loading a YAML \texttt{DefinitionDocument}, and the engine works
out the operations needed to make the tree match the document. The
load operation goes through the five stages of \Cref{fig:load-flow}:

\begin{enumerate}
  \item \textbf{Parse.} The YAML is read with \texttt{ruamel.yaml}
        and validated by Pydantic against the document schema.
        Syntax errors are caught here.
  \item \textbf{Diff.} The engine compares the document against the
        current tree. The result is a list of intended operations
        (create this folder, add this descriptor, update this
        template) without yet touching anything.
  \item \textbf{Validate.} Every intended operation is validated:
        does the user have permission, is the hypervisor name
        resolved against an existing definition, is the target
        datastore present? Validation \emph{rejects} the load when
        any check fails, \emph{before} any job is created.
  \item \textbf{Plan.} The same flow with execute disabled. Useful
        as a dry-run command. Returns the intended operations and
        the validation result without touching the tree.
  \item \textbf{Execute.} The operations are translated into
        Commands and submitted to the job engine. Each job is
        persisted, executed and logged.
\end{enumerate}

The strict separation between validation and execution is a
deliberate choice (DEC-032). A failed validation never creates a
job; this keeps the job log free of records that exist only to mark
``user submitted a bad request.'' Inside execution, however, a job
that fails after passing validation \emph{does} create a record,
because that is exactly the kind of incident the operator needs to
see.

\begin{figure}[ht]
  \centering
  \includegraphics[width=0.75\linewidth]{load-flow.pdf}
  \caption{Stages of \texttt{load} / \texttt{plan}: validation
    rejects \emph{before} any state change.}
  \label{fig:load-flow}
\end{figure}

\subsection{Descriptor state machine (DEC-022, DEC-032)}
\label{subsec:state-machine}

Every descriptor has a state, tracked by the orchestrator
independently of the runtime state of the underlying VM. There are
four states: \texttt{provisioned} (the descriptor exists but no VM
has been created in the hypervisor yet), \texttt{deployed} (the
descriptor matches a real VM in the hypervisor),
\texttt{unreachable} (the hypervisor cannot be contacted; the actual
state of the VM is unknown) and \texttt{broken} (an operation
failed in a way that leaves the orchestrator and the hypervisor
out of sync). The transitions are shown in \Cref{fig:state-machine}.

Two characteristics of this state machine deserve emphasis. First,
\emph{state changes only as the result of a completed job}: there is
no background polling that silently transitions a descriptor. This
makes every state change auditable through the job log. Second, the
runtime state of the VM (running, stopped, paused) is reported by
the hypervisor on demand and not persisted by the orchestrator. This
keeps the orchestrator's model small and prevents drift between two
state representations that would inevitably diverge.

The \texttt{broken} state has the precise definition added in
DEC-032: a descriptor is \texttt{broken} only after a job that was
expected to converge the descriptor with the hypervisor (typically a
deploy or undeploy) failed in a way that left a VM partially
created or partially destroyed. The recovery path is the
administrator's \texttt{force-undeploy} operation: clean up whatever
exists in the hypervisor and return the descriptor to
\texttt{provisioned}.

\begin{figure}[ht]
  \centering
  \includegraphics[width=0.85\linewidth]{state-machine.pdf}
  \caption{Descriptor state machine. Runtime states (running,
    stopped, paused) belong to the hypervisor and are not persisted by
    the orchestrator.}
  \label{fig:state-machine}
\end{figure}

\subsection{REST as the public boundary}
\label{subsec:rest-boundary}

The decision to expose a REST API as the boundary of the daemon was
made in Sprint~3 and deserves an explicit defence because it shapes
the entire client architecture. The argument for REST is \emph{not}
that the CLI needs to be usable remotely without SSH; the
administrator already has SSH. The argument is that the REST API
exists as a \emph{public boundary for external integrations}.
External clients (CI/CD pipelines, GitOps tools, future teaching
applications, third-party scripts) have a clean and documented
contract to depend on. The CLI of UnivOrch then \emph{also} uses
that contract, which removes any temptation to add a back door for
local use and keeps the in-process and out-of-process paths
identical.

An intermediate design (the CLI talks to the engine in-process when
running on the same host and over HTTP when remote) was considered
and discarded during Sprint~3. The advantage of fewer moving parts
was real, but the cost in duplicated logic and divergent behaviour
between the two paths was higher. The decision aligns UnivOrch with
mature peers: \texttt{docker}, \texttt{kubectl} and \texttt{podman}
all run as thin HTTP clients of a daemon, even when both live on
the same machine.

%==================================================================
\section{Implementation}
\label{sec:implementacion}

This section presents the implementation by layer rather than by
sprint. The sprints are documented in \Cref{ch:metodologia}; for the
description of the product, the layered view is clearer. The code
lives under \texttt{src/univorch/}; each subsection below
corresponds to one subdirectory of that tree.

\subsection{Orchestrator core}

The core consists of the Pydantic models, the resolver, the job
engine, the repositories and the facade service.

\textbf{Models}. The domain types are implemented as Pydantic
\texttt{BaseModel} classes with strict configuration
(\texttt{extra="forbid"}, field validators for the path syntax). The
two families discussed in \Cref{subsec:domain-model} are kept in
separate modules to avoid accidental cross-contamination between
the persisted form and the document form.

\textbf{Resolver}. The resolver is a pure function: given a folder
path and the current tree, it returns the effective definition of
that folder, with cascade inheritance applied and the lexical
closure resolved for any imported template. ``Pure'' has a precise
meaning here: no side effects, no I/O, no exceptions other than for
truly invalid input. This makes the resolver the ideal candidate for
property-based testing with Hypothesis (\Cref{sec:pruebas}). The
inheritance rule is encoded once per data type
(scalar/list/dictionary) and the same rule applies to every field.
The walker that follows the import chain through ancestor folders
is generic, parameterized by the resource type (template,
hypervisor); adding a new inheritable resource type in the future is
a one-line change.

\textbf{Job engine and Commands}. The job engine accepts a Command
and turns it into a Job. The Command is responsible for
\emph{validation} (a single method that returns a list of
human-readable error messages and does not touch the database) and
\emph{execution} (a single method that mutates the world). The
engine wraps the Command, creates a Job record in TinyDB,
transitions the descriptor state if needed, captures exceptions and
records the outcome. Six Commands are implemented in v1:
\texttt{DeployCommand}, \texttt{UndeployCommand},
\texttt{StartCommand}, \texttt{StopCommand},
\texttt{CreateFolderCommand} and \texttt{CreateDescriptorCommand}.

\textbf{Repositories}. There is one repository per aggregate root
(folder, descriptor, job). The interface is small: \texttt{save},
\texttt{get}, \texttt{exists}, \texttt{subtree} or \texttt{find\_by\_status},
\texttt{delete}. The TinyDB implementation operates on Pydantic
models through \texttt{model\_dump()} and \texttt{model\_validate()}.
A different backend (MongoDB, PostgreSQL) would replace the body of
these methods without changing any caller.

\textbf{Facade service}. The \texttt{OrchestratorService} class is
the public face of the engine. It exposes one method per public
operation (\texttt{deploy}, \texttt{undeploy}, \texttt{start},
\texttt{stop}, \texttt{status}, \texttt{list\_tree}, \texttt{inspect},
\texttt{load}, \texttt{folder\_exists}) and resolves the connector
session through the connection pool. The facade is the only place
where coordination between repositories, the resolver, the job
engine and the connector pool happens.

\subsection{Mock connector}

The mock connector is the only implementation of the
\texttt{HypervisorConnector} abstract base class delivered in v1.
It keeps its state in memory (a dictionary of fake VMs, a counter
of created VM identifiers, a list of pre-loaded demo templates) and
implements every method of the contract with the same observable
behaviour as a real hypervisor. The mock is not a stub: it actually
clones the requested template into a new VM record, transitions
its runtime state under \texttt{start} and \texttt{stop} and reports
the right value under \texttt{get\_status}. The whole system runs
end-to-end against this connector, without any real hypervisor;
this is what makes the demo of \Cref{sec:demo} reproducible on any
laptop.

\subsection{CLI interface}

The CLI is implemented with \texttt{cmd2}, which provides both a
batch mode (one command at a time, suitable for shell scripts) and
an interactive REPL (with history, tab completion and built-in
\texttt{help}). Argument parsing for every command uses the
\texttt{argparse} integration of \texttt{cmd2}; this gives
auto-generated, coloured help for free.

The CLI does not embed the engine. It accepts an instance of
\texttt{OrchestratorAPI} at construction time and forwards every
command to it. By default the instance is an
\texttt{HttpServiceClient} configured against the daemon's URL; a
test that wants to drive the CLI without a daemon passes an
in-process \texttt{OrchestratorService} instead. The
\texttt{HttpServiceClient} translates HTTP transport errors
(connection refused, timeout, 4xx response) into
\texttt{OperationError} exceptions with actionable messages, so the
user does not see HTTP-level error codes in the terminal.

\subsection{REST daemon}

The daemon is a FastAPI application served by Uvicorn. The
application is built by \texttt{create\_app(service)}, which takes
an instance of the facade and registers one endpoint per public
operation. Paths are passed as query parameters
(\texttt{?path=/lab/student01}) so they remain easy to read in
URLs and easy to test from the command line with \texttt{curl}.
Error codes follow a small and explicit convention: 200~OK for
successful operations, 400~Bad Request for failed validation
(returning the list of error messages), 422~Unprocessable Entity for
malformed inputs that fail Pydantic validation before reaching the
service. There is a \texttt{GET /api/v1/health} endpoint for
liveness checks.

The daemon entry point lives in
\texttt{src/univorch/interfaces/rest/\_\_main\_\_.py} and is
registered in \texttt{pyproject.toml} as the \texttt{univorchd}
console script. It composes a real \texttt{OrchestratorService} (with
TinyDB repositories on the path declared by the environment
variable \texttt{UNIVORCH\_DB\_PATH}) and runs Uvicorn on the port
declared by \texttt{UNIVORCH\_PORT}.

\subsection{Web interface (read-only)}

A first version of the web interface is delivered in Sprint~4 as a
read-only view of the tree. It is built with NiceGUI and mounted on
the same Uvicorn process as the REST application via
\texttt{ui.run\_with(fastapi\_app, \dots)}: a single port serves both
\texttt{/api/v1/*} and the web pages. This is a deliberate departure
from the Sprint~3 ``every client is HTTP'' rule: the web is part of
the daemon, not an external HTTP client of it. The external clients
(CLI, the teaching subcommands, third-party integrations) keep
going through the REST API.

The v1 has three pages. The index page shows a header, four summary
cards (\emph{total VMs}, \emph{provisioned}, \emph{deployed},
\emph{broken or unreachable}) and an expandable tree of folders and
descriptors. The descriptor nodes carry an icon and a badge that
match the state vocabulary of the CLI (grey for
\texttt{provisioned}, green for \texttt{deployed}, red for
\texttt{broken}, yellow for \texttt{unreachable}). Clicking a
descriptor node navigates to the descriptor detail page, which
calls \texttt{service.inspect(path, resolved=True)} and
\texttt{service.status(path)} and renders the eleven relevant
fields of the effective definition. Folder nodes only expand or
collapse, with no navigation.

Write operations (deploy, start, stop, undeploy, load),
authentication and the workstation/desk view for students remain on
the roadmap; the read-only v1 is the proof that the boundary with
the orchestrator service works correctly through a graphical client.

\subsection{Teaching application (layer 2)}

The teaching application is the only layer-2 application implemented,
and it is the practical demonstration that the two-layer split
(DEC-004) works: a complete domain workflow built entirely on the
generic engine, without the engine learning any teaching vocabulary.
It lives in \texttt{src/univorch/teaching/} and is exposed as a
\texttt{teach} subcommand group of the CLI; it talks to the core
through the same \texttt{OrchestratorAPI} as every other client.

The only change it required in the core is a single opaque
\texttt{metadata} dictionary on the folder model (DEC-038). The core
stores it and never interprets it; the teaching layer puts
\texttt{\{kind: subject, desktop: [\dots]\}} there to mark a folder as
a subject and to record which templates make up a student's set of
machines. A different layer-2 application would reuse the same field
for its own markers, so the core stays domain-agnostic.

The workflow has two inputs and a handful of commands:

\begin{description}
  \item[\texttt{teach load-subject}] validates that a YAML document is
        a well-formed subject (marked \texttt{kind: subject},
        non-empty desktop, every desktop template resolvable, no child
        folders) and, only then, loads it with the core
        \texttt{load}. A malformed subject, or the wrong file
        altogether, is rejected with a single actionable message
        before anything is created.
  \item[\texttt{teach load-students}] reads the subject's desktop from
        its metadata and generates, under the subject, one folder per
        student (named after the username) and one descriptor per
        desktop template inside each. Each descriptor carries
        \texttt{use template: <name>} and the student folder carries
        \texttt{import: *}, so the descriptors resolve the subject's
        templates by cascade --- the generated tree is a normal core
        document the resolver already understands. The operation is
        add-only: existing students are no-ops, none are removed
        (DEC-039).
  \item[\texttt{teach save-students}] exports the current roster, the
        inverse of \texttt{load-students} at the list level.
\end{description}

Two design choices keep the layer thin. First, there is no new
\texttt{use desktop:} keyword: because the application creates the
descriptors, it writes \texttt{use template: <name>} directly, which
the core already resolves. Second, the operations are decoupled from
the CLI into a \texttt{teaching/operations.py} module, so they are
tested against a real in-process service without a daemon. An
end-to-end demonstration --- load a subject of six machines, load a
roster of ten students (one of them an email address), and deploy one
of the resulting machines --- runs against the mock connector and
produces sixty provisioned descriptors that each resolve their
template by inheritance.

Deploying a whole subject in one batch command, removing students
that left the roster, email notifications and a dedicated teaching
web view are designed and left as future work.

\subsection{Packaging and distribution}

The production image is built from a multi-stage Dockerfile. The
first stage uses \texttt{uv} to resolve and install only the
runtime dependencies. The second stage is a slim Python 3.12 base
that copies the resolved virtual environment, the application code
and the entry points. The resulting image is around 200~MB and
takes a few seconds to pull. The image is published to
\texttt{ghcr.io/clamaveruma/univorch} by the \texttt{publish.yml}
workflow, triggered by signed tags of the form \texttt{vX.Y.Z}.

A \texttt{docker-compose.yml} file declares the service, the
named volume (\texttt{univorch\_data}) where the TinyDB file
lives, the port mapping and the environment variables. A thin
shell wrapper, \texttt{univorch.sh}, exposes a convenient command
surface (\texttt{start}, \texttt{stop}, \texttt{restart},
\texttt{status}, \texttt{logs}, \texttt{cli}).

For new users, an \texttt{install.sh} script is published from a
fixed URL and does the whole installation in one command. The
script detects the local Docker version, picks a free port if 8080
is taken, falls back from \texttt{docker compose} to
\texttt{docker-compose} v1 if needed and explains what it is going
to do before doing it. The script is documented in the supervisor's
tutorial (\texttt{docs/tutorial-profesor.md}).

%==================================================================
\section{Testing strategy}
\label{sec:pruebas}

The testing strategy is a small but consistent application of the
test pyramid. Most tests are unit tests; a strong layer of
integration tests covers the facade and the REST endpoints; a
minority of tests use property-based testing to exercise the
resolver. There are no end-to-end tests against a real hypervisor
in v1, because no real connector exists yet; the demo of
\Cref{sec:demo} is the closest analogue.

\subsection{Unit tests}

Unit tests target a single class or module. They are fast (the full
unit suite runs in well under one second) and they pin behaviour at
the granularity needed to refactor with confidence. The Pydantic
models are tested for round-trip serialization and for field
validation. The Commands are tested with both a successful and a
failing case for each branch of their validation. The mock
connector is tested by exercising every method of the
\texttt{HypervisorConnector} contract.

\subsection{Integration tests}

Integration tests target the facade and exercise it through the
complete chain: facade $\to$ resolver $\to$ job engine $\to$
repositories $\to$ TinyDB $\to$ mock connector. The repositories
are configured against an in-memory TinyDB so the tests stay fast
and isolated from disk; the mock connector is used in its
\texttt{empty()} variant so every test starts from a known state.
The integration tests are the ones that catch interaction bugs
between pieces that look correct in isolation.

\subsection{Property-based tests}

The resolver is a pure function over plain data. Property-based
testing with Hypothesis is well suited to it.
The strategy generates random trees of folders and descriptors with
random imports, templates and hypervisor definitions, and checks
properties that should hold for any input:

\begin{itemize}
  \item Idempotence: resolving the same descriptor twice returns
        the same result.
  \item Closure: a template imported from $A$ resolves its internal
        references from $A$, regardless of the folder where the
        import happens.
  \item Bounded import: a folder cannot see a resource that no
        ancestor explicitly imported.
\end{itemize}

Hypothesis shrinks failing cases to a minimal counterexample
automatically; in practice this has caught two subtle bugs in
the inheritance resolution during the development of Sprint~2.

\subsection{REST tests}

The REST endpoints are tested with the \texttt{httpx.AsyncClient}
of FastAPI, which routes requests directly to the ASGI application
in-process. This catches client/server divergences (a missing field
in the request schema, a wrong status code) without spinning up a
real Uvicorn server. The same tests also exercise the
\texttt{HttpServiceClient} of the CLI: the client is constructed
against the in-process test transport, which keeps the test fast
and deterministic. This is one of the cleanest benefits of the
Sprint~3.2 decision to make the CLI a pure HTTP client: the same
test suite covers the wire format and the client code at once.

\subsection{Coverage and metrics}

\Cref{tab:coverage} reports the line and branch coverage at the
close of Sprint~3, measured by \texttt{pytest-cov} on the central
modules. Coverage is above 95\% in every module that contains
domain logic; the lower number on the CLI module reflects mostly
interactive code paths and the launch wiring that the test suite
intentionally does not exercise.

\begin{table}[ht]
  \centering
  \small
  \begin{tabular}{l r r}
    \toprule
    \textbf{Module} & \textbf{Line coverage} & \textbf{Branch coverage} \\
    \midrule
    \texttt{models.py}            & 100\% & 100\% \\
    \texttt{resolver.py}          & 100\% & 100\% \\
    \texttt{service.py}           & 97\%  & 95\% \\
    \texttt{jobs/commands.py}     & 100\% & 100\% \\
    \texttt{jobs/engine.py}       & 100\% & 100\% \\
    \texttt{persistence/tinydb/}  & 100\% & 100\% \\
    \texttt{connectors/mock.py}   & 100\% & 100\% \\
    \texttt{interfaces/cli/}      & 92\%  & 88\% \\
    \texttt{interfaces/rest/}     & 95\%  & 90\% \\
    \bottomrule
  \end{tabular}
  \caption{Coverage on the core modules (286 tests passing, including
    the teaching layer).}
  \label{tab:coverage}
\end{table}

%==================================================================
\section{Functional demonstration}
\label{sec:demo}

This section describes the end-to-end demonstration that runs
against the mock connector, without any real hypervisor. The same
demonstration is in \texttt{demo/README.md} of the repository, with
copy-pasteable commands. Screenshots will be added before the final
submission of the report; provisionally the demonstration is
described in prose.

\subsection{Starting the daemon}

The administrator pulls the public image and starts the service in
one command:

\begin{verbatim}
docker run -d --name univorch -p 8080:8080 \
       -v univorch_data:/data \
       ghcr.io/clamaveruma/univorch:latest
\end{verbatim}

A liveness check at \texttt{http://localhost:8080/api/v1/health}
returns \texttt{\{"status":"ok"\}}. The OpenAPI specification is
auto-generated by FastAPI and available at
\texttt{/api/v1/openapi.json}; the Swagger UI is at
\texttt{/api/v1/docs} and is the easiest way to explore the API
without writing any code.

\subsection{Loading the demo tree}

The demo tree is a small YAML file describing one administrative
folder \texttt{/lab} (with the hypervisor definition and a base
template), one folder \texttt{/lab/networks} (the subject) and
three descriptor entries \texttt{student01}, \texttt{student02},
\texttt{student03} that use the imported template:

\begin{verbatim}
kind: definition
version: "1"
folders:
  lab:
    description: "Computer Science Lab"
    define hypervisors:
      mock01: { type: mock }
    define templates:
      linux-vm:
        use hypervisor: mock01
        base_vm: linux-base
    folders:
      networks:
        description: "Computer Networks course"
        import: [linux-vm]
        descriptors:
          student01: { use template: linux-vm, cpu: 2 }
          student02: { use template: linux-vm, cpu: 2 }
          student03: { use template: linux-vm, cpu: 4 }
\end{verbatim}

The administrator loads the document into the orchestrator:

\begin{verbatim}
univorch load demo/setup.yml /
\end{verbatim}

The CLI prints the list of operations performed, one line per
folder or descriptor created. After this command, \texttt{univorch
tree /} renders the descriptor tree with state glyphs next to each
node: \texttt{[ ]} for \texttt{provisioned}, \texttt{[*]} for
\texttt{deployed}, in red for \texttt{broken}, in yellow for
\texttt{unreachable}.

\subsection{Inspecting the inheritance}

The \texttt{inspect} command shows the effective definition of any
folder or descriptor with the inheritance applied. Three modes are
supported. \texttt{--local} prints only what is written at the node
itself, with no merging. \texttt{--expanded} prints the result of
applying imports and cascade inheritance but without following
references inside templates. The default mode prints the fully
resolved definition: the hypervisor name, the base VM, the hardware
parameters and any local overrides.

This is how the demonstration shows the lexical closure: the
descriptor \texttt{student03} declares no hypervisor and no base VM
explicitly; \texttt{inspect} reveals that it inherits the template
\texttt{linux-vm} (defined in \texttt{/lab}, not in
\texttt{/lab/networks}) and that the template itself resolves
\texttt{use hypervisor: mock01} against the same folder
\texttt{/lab} where it was defined. The override of \texttt{cpu} to
4 appears explicitly as a local property.

\subsection{Full life cycle}

The administrator runs the complete life cycle on one of the
descriptors:

\begin{verbatim}
univorch deploy /lab/networks/student01
univorch start  /lab/networks/student01
univorch status /lab/networks/student01
univorch stop   /lab/networks/student01
univorch undeploy /lab/networks/student01
\end{verbatim}

Each command produces a Job that is persisted in the daemon. The
\texttt{deploy} job transitions the descriptor from
\texttt{provisioned} to \texttt{deployed}; the mock connector
records the creation of a new VM with a fresh identifier. The
\texttt{start} job transitions the runtime state of the VM but
does not change the descriptor state. The \texttt{status} command
reads both states from the connector. The \texttt{stop} and
\texttt{undeploy} jobs reverse the sequence; after
\texttt{undeploy}, the descriptor is back in \texttt{provisioned}
and the mock connector has no record of the VM.

\subsection{Exercising the REST API}

The same operations can be invoked directly through the REST API
with \texttt{curl}:

\begin{verbatim}
curl -X POST 'http://localhost:8080/api/v1/deploy?path=/lab/networks/student01'
curl 'http://localhost:8080/api/v1/status?path=/lab/networks/student01'
\end{verbatim}

This is the same code path the CLI uses, and it is also the entry
point that future external integrations (a CI pipeline, a teaching
front-end, a monitoring agent) would use. The Swagger UI displays
every endpoint with its parameter schema, response schema and
example values; an external developer who has never seen UnivOrch
can navigate it in a few minutes.

%==================================================================
\section{Comparison with esxobjects and yamlinfr}
\label{sec:comparativa}

The two libraries written by the supervisor before this work,
\texttt{esxobjects} and \texttt{yamlinfr}, are the operational
precedent of UnivOrch. This section compares them with UnivOrch on
the same axes that matter for a proof of concept aimed at
maintainability and reuse. \Cref{tab:comparison} summarizes the
comparison.

\begin{table}[ht]
  \centering
  \small
  \begin{tabular}{p{0.20\linewidth} p{0.23\linewidth} p{0.23\linewidth} p{0.25\linewidth}}
    \toprule
    \textbf{Axis} & \textbf{\texttt{esxobjects}} & \textbf{\texttt{yamlinfr}} & \textbf{UnivOrch} \\
    \midrule
    Scope &
    VMware vSphere API wrapper. &
    YAML-to-deployment script over \texttt{esxobjects} + external
    IPAM. &
    Long-running orchestration service, hypervisor-agnostic. \\
    Hypervisor coupling &
    VMware only. &
    VMware only (via \texttt{esxobjects}). &
    Pluggable: mock today, VMware and Proxmox planned. \\
    Model &
    Object-oriented wrapper around vSphere objects. &
    Flat YAML, one infrastructure per file. &
    Tree of folders with cascade inheritance and lexical closure of
    templates. \\
    Persistence &
    None: each invocation is independent. &
    None: stateless deployment. &
    TinyDB document database (replaceable by MongoDB). \\
    Public boundary &
    Python library. &
    Command-line script. &
    REST API with OpenAPI specification. \\
    Test coverage &
    Limited. &
    Limited. &
    286 tests, $>$95\% coverage on the core. \\
    Distribution &
    Internal department use. &
    Internal department use. &
    Public Docker image in \texttt{ghcr.io}, single-command install. \\
    \bottomrule
  \end{tabular}
  \caption{Comparison of \texttt{esxobjects}, \texttt{yamlinfr} and
    UnivOrch.}
  \label{tab:comparison}
\end{table}

The comparison must be read honestly. The two previous libraries
\emph{solved their problem} -- a real working deployment system
that the supervisor and his students have used for actual courses.
UnivOrch is, today, a proof of concept of a more ambitious model:
it generalizes the hypervisor coupling, lifts the deployment script
into a long-running service, formalizes the data model into a tree
with inheritance, exposes a public REST contract and adds the
infrastructure (test discipline, CI/CD, packaged distribution) that
the previous libraries did not have time to develop. It does \emph{not
yet} replace them operationally: the real VMware connector is still
to be implemented, so an actual deployment against the teaching lab
cannot run on UnivOrch today.

The intended path forward is to write the VMware connector against
the abstract interface defined in Sprint~1, validate it against the
same lab that \texttt{yamlinfr} drives today, and then deprecate
\texttt{yamlinfr} in favour of UnivOrch with the teaching
application of Sprint~5 on top. This work is outside the scope of
this TFG, but the present proof of concept proves that the model
underneath is sound.
